% PiLENS v2 manuscript draft (roadmap Step 11).
% Written with IEEEtran so it compiles anywhere (Overleaf: New Project -> Upload).
% For submission, paste the sections into the official IEEE Access template.
% Every \TBD{} is a number that must come from docs/v2_results.md (measured only).
\documentclass[journal]{IEEEtran}
\usepackage{amsmath,amssymb,graphicx,booktabs,xcolor,cite,url}
\newcommand{\TBD}[1]{\textcolor{red}{[TBD: #1]}}

\begin{document}

\title{PiLENS: A Two-Stage Video Anomaly Detection Cascade for Real-Time\\
Night-Vision Surveillance on the Raspberry Pi 5}

\author{Huzaifa~Asad, \TBD{co-authors and affiliation, order to be confirmed with supervisor}}

\maketitle

\begin{abstract}
Automated surveillance on low-cost edge hardware must detect suspicious activity in real time,
at night, and without a GPU. We present PiLENS, a Raspberry Pi~5 system with a NoIR camera and
infrared illumination that runs a two-stage cascade on a single shared X3D-S backbone. Stage~1,
a weakly supervised multiple-instance-learning (MIL) head, scores every 2.4\,s clip as normal
or suspicious; stage~2, a 14-class UCF-Crime head, runs only when a $k$-of-$n$ temporal vote
raises an alert and reuses cached clip features, so it adds almost no compute. On the official
UCF-Crime anomaly split the binary stage reaches a video-level AUC of \TBD{video AUC $\pm$ std}
and a frame-level AUC of \TBD{frame AUC $\pm$ std}; 14-class recognition on the official four
folds reaches \TBD{accuracy}\% (macro-F1 \TBD{macro-F1}\%), which we treat as assistive. Exported
to ONNX Runtime, the cascade processes one clip in \TBD{median clip latency}\,ms (p95
\TBD{p95}) on the Pi~5 CPU while the camera pipeline sustains \TBD{stream fps}\,frames/s, with
an event-to-alert time of \TBD{median}\,s and \TBD{FA/h} false alarms per hour. All latency,
throughput and temperature figures are measured on the device. Code, split files and the
evaluation protocol are released.
\end{abstract}

\begin{IEEEkeywords}
Video anomaly detection, edge computing, Raspberry Pi, night vision, multiple instance
learning, X3D, ONNX Runtime, surveillance.
\end{IEEEkeywords}

\section{Introduction}
CCTV in universities, warehouses and residential areas is still watched by people. Manual
monitoring of many feeds is tiring and error prone, and footage is usually reviewed only after
an incident. Night-time scenes make this worse: standard cameras lose detail, and operator
attention is lowest.

Deep video models can flag abnormal events automatically~\cite{sultani2018}, but most published
systems assume a desktop GPU or a cloud server, report accuracy only, and are evaluated on
daylight footage. A deployable system also has to answer engineering questions that accuracy
tables do not: how long one decision takes on the target CPU, how many frames per second the
whole pipeline sustains, how quickly an alert follows the start of an event, how many false
alarms an operator receives per hour, and whether the device throttles under sustained load.

This paper presents PiLENS, an end-to-end night-vision surveillance system on a Raspberry Pi~5.
Our contributions are:
\begin{enumerate}
\item A two-stage cascade on one shared backbone: a binary MIL anomaly head on every clip and a
14-class head only on alerts, operating on cached features (Section~\ref{sec:method}).
\item A measured edge benchmark on the Pi~5: per-stage clip latency, stream throughput,
event-to-alert time, false alarms per hour, CPU/RAM use and temperature with and without active
cooling (Section~\ref{sec:results}).
\item Night-vision handling: grayscale-to-three-channel preprocessing shared by training and
inference, fixed camera exposure, a motion gate robust to IR-attracted insects, and a day/night
split of the test set to expose a ``darkness = anomaly'' shortcut.
\item A reproducible protocol on the official UCF-Crime splits with video-level validation,
duplicate removal and tuning on validation only, released with the code.
\end{enumerate}

\section{Related Work}
\subsection{Weakly supervised video anomaly detection}
Sultani et al.~\cite{sultani2018} introduced UCF-Crime and the MIL ranking formulation, in which
each untrimmed video is a bag of temporal segments with only a video-level label. Later work
improved feature magnitude learning~\cite{tian2021rtfm} and temporal context. These methods
typically use heavy I3D~\cite{carreira2017} features on GPUs. \TBD{add 3--5 recent VAD papers}

\subsection{Efficient video recognition}
X3D~\cite{feichtenhofer2020x3d} expands a tiny 2D network along temporal, spatial and width axes;
X3D-S needs about 2\,GFLOPs per clip. MoViNets~\cite{kondratyuk2021movinets} add causal streaming
buffers for online inference. Both are pretrained on Kinetics~\cite{kay2017kinetics}.

\subsection{Edge surveillance systems}
\TBD{Raspberry Pi / Jetson surveillance and anomaly detection systems; note which report
measured on-device latency, throughput and temperature, and which report accuracy only.}

\section{System Design}\label{sec:system}
\subsection{Hardware}
A Raspberry Pi~5 (8\,GB) runs everything on its CPU. A NoIR camera module with an infrared
illuminator provides near-grayscale night frames; exposure and gain are fixed to avoid
auto-exposure flicker. A GPIO LED and buzzer give local alerts; alert clips are e-mailed; an
MJPEG stream is reachable only through a Tailscale private network.
\TBD{Fig.: hardware photo and circuit diagram (assets/).}

\subsection{Real-time pipeline}
The camera delivers 30 frames/s into a ring buffer (capture thread, never blocked). A motion
gate runs MOG2 background subtraction~\cite{zivkovic2004} on a blurred $320\times180$ grayscale
frame; a blob counts only if its area exceeds a minimum fraction of the frame and it persists
over consecutive checks, which suppresses insects near the IR LEDs and light flicker. Every hop
($h=0.6$\,s), if the gate is open, the inference thread takes the latest 13-frame clip (stride
0.2\,s, 2.4\,s span) from the buffer, extracts one 2048-d X3D-S feature, and scores it with the
binary head. An alert is raised when $k$ of the last $n$ clips exceed the validation threshold
($k{=}2$, $n{=}3$), followed by a cooldown. Stage~2 then averages the cached features of the last
3--5 clips and returns the top-3 classes; below a confidence floor the label is ``Unknown
suspicious''. A separate alert thread drives the GPIO, writes the preceding 5\,s to disk and
sends the e-mail, so video encoding and networking never stall inference.
\TBD{Fig.: block diagram of the threads and the cascade.}

\section{Method}\label{sec:method}
\subsection{Clip features}
Each untrimmed training video is divided into $T{=}32$ equal temporal segments, and one clip of
13 frames (stride 6 at 30\,fps) is sampled at the centre of each segment. Frames are resized to
a short side of 182 pixels and centre-cropped to $160\times160$. The X3D-S backbone (Kinetics-400
weights, classifier removed) maps each clip to a 2048-d pre-logit vector. The backbone is frozen:
end-to-end fine-tuning on trimmed clips overfit (train $\approx$92\%, validation 27--34\%).

\subsection{Stage 1: binary MIL head}
Let $s^a_t$ and $s^n_t$ be the clip scores of an anomalous and a normal video. With $\mathrm{top}_k$
the mean of the $k{=}3$ highest scores, the loss is
\begin{equation}
\mathcal{L} = \max\!\big(0,\,1-\mathrm{top}_k(s^a)+\mathrm{top}_k(s^n)\big)
+ \lambda_1 \sum_{t=1}^{T-1}(s^a_t-s^a_{t+1})^2 + \lambda_2 \sum_{t=1}^{T} s^a_t,
\end{equation}
with $\lambda_1=\lambda_2=8\times10^{-5}$. We compare a per-clip MLP ($2048{\to}512{\to}32{\to}1$)
and a temporal convolution head. Each step samples equal numbers of anomalous and normal bags.

\subsection{Stage 2: 14-class head}
On the same cached features we train a linear head on the mean clip feature, a top-$k$ MIL head
and a temporal-convolution MIL head, with class-balanced sampling and label smoothing.

\subsection{Deployment}
Backbone and heads are exported separately to ONNX (opset 17); outputs match PyTorch to
\TBD{max |diff|} on real clips. Inference uses ONNX Runtime on the CPU with
\TBD{threads} intra-op threads. \TBD{INT8: static QDQ quantization calibrated on day and night clips.}

\section{Dataset and Protocol}\label{sec:data}
We use UCF-Crime~\cite{sultani2018} (1950 files, 320$\times$240, 30\,fps). Fifty Normal videos
appear twice under different folders; byte-identical copies are removed, leaving 1900 unique
videos. We follow the official splits: the 14-class 4-fold split (38 training and 12 test videos
per class) and the anomaly split (1610 training, 290 test). From each training list, 15\% of the
videos are held out per class as validation. Every split is at video level, so all clips of a
video belong to one split. Epochs, checkpoints and the alert threshold are selected on validation
only; each test set is evaluated once with fixed settings. Frame-level ground truth uses each
test video's actual frame count.

Roughly half of the anomaly videos are recorded at night. Normal videos carry no day/night
label, so we infer it from frame saturation and brightness with thresholds calibrated on 947
manually labelled anomaly videos (balanced accuracy \TBD{value}), and report day and night
test subsets separately.

\section{Experimental Setup}
Features and heads are trained on Kaggle (2$\times$NVIDIA T4). Results are mean $\pm$ standard
deviation over folds (Task~A, 3 seeds per fold) or over 5 seeds (Task~B). Metrics: accuracy,
macro-F1, per-class precision/recall/F1 and confusion matrix (Task~A); video AUC, frame AUC,
frame AP and precision/recall at the validation threshold (Task~B). The edge benchmark uses
\TBD{Pi OS version, Python, ONNX Runtime version}; each latency figure follows 20 warm-up and
200 timed runs and reports median and 95th percentile. Thermal runs last 30 minutes with and
without active cooling.

\section{Results}\label{sec:results}
\TBD{Copy Tables 1--5 from docs/v2_results.md. Preliminary Kaggle numbers in the notes:
14-class linear 33.2$\pm$0.3\% acc / 31.8$\pm$0.3\% macro-F1 (chance 7.1\%);
binary MLP video AUC 0.912$\pm$0.005, frame AUC 0.806$\pm$0.010, frame AP 0.205.
Use only numbers reproduced in Step 3.}

\subsection{Recognition and anomaly detection}
\subsection{On-device latency and throughput}
\subsection{Alert behaviour: event-to-alert time and false alarms}
\subsection{Ablations}
v1 (ResNet18-LSTM) vs.\ X3D-S vs.\ MoViNet; crop strategy; RGB vs.\ grayscale-3ch; motion gate
and $k$-of-$n$ rule on/off; FP32 vs.\ INT8; day vs.\ night.

\section{Discussion and Limitations}
The binary stage is strong while 14-class recognition is weak, so the cascade is built on the
binary decision and the class label is presented as assistive with a confidence floor. UCF-Crime
is 320$\times$240 CCTV footage, unlike the Pi NoIR camera; we therefore evaluate on our own
recorded clips \TBD{external test}. If Normal videos are mostly daytime, a model can learn
darkness as a cue for anomaly; the day/night breakdown tests this. Sustained load can throttle
the CPU; we report temperature and throughput over time.

\section{Ethics}
Surveillance footage contains identifiable people. We use UCF-Crime under its terms of use, and
our own recordings were made with written consent of all participants
\TBD{ethics approval reference}. Alert e-mails and remote access are restricted to a private
Tailscale network; credentials are kept out of the code.

\section{Conclusion}
\TBD{2--3 sentences after results are final.}

\bibliographystyle{IEEEtran}
\bibliography{refs}
\end{document}
